\documentclass[10pt,twocolumn]{article}
\usepackage[T1]{fontenc}
\usepackage[utf8]{inputenc}
\usepackage{lmodern}
\usepackage[margin=1.8cm,top=2.4cm,bottom=2cm,columnsep=0.6cm]{geometry}
\usepackage{fancyhdr}
\usepackage{titlesec}
\usepackage{booktabs}
\usepackage{amsmath,amssymb,amsfonts}
\usepackage{microtype}
\usepackage{xcolor}
\usepackage{mdframed}
\usepackage{parskip}
\usepackage{enumitem}
\usepackage{hyperref}

\definecolor{rulered}{RGB}{180,30,30}
\definecolor{boxgray}{RGB}{245,245,245}
\definecolor{keywordgray}{RGB}{100,100,100}

\pagestyle{fancy}
\fancyhf{}
\fancyhead[L]{\textsc{\small Claude Mag}}
\fancyhead[C]{\small\textit{Machine Learning Research Digest}}
\fancyhead[R]{\small 2026-11-18}
\fancyfoot[C]{\small\thepage}
\renewcommand{\headrulewidth}{0.8pt}

\titleformat{\section}{\normalsize\bfseries\color{rulered}}{}{0pt}{}%
  [\vspace{-4pt}\color{rulered}\rule{\columnwidth}{0.4pt}]
\titlespacing{\section}{0pt}{8pt}{4pt}

\hypersetup{colorlinks=true,urlcolor=rulered,linkcolor=black}

\begin{document}
\setlength{\parindent}{0pt}
\setlength{\parskip}{4pt}

{\small\color{keywordgray}\textbf{Keywords:}
  Sim-to-Real \textbullet{} Robot Learning \textbullet{} Reinforcement Learning \textbullet{} Reasoning Models}\par\vspace{4pt}

{\LARGE\bfseries\raggedright Training Reasoning Models in Simulation to Natively Use Real Robots}\par\vspace{4pt}
{\large\itshape\raggedright PEAR Trains a 27B Model Entirely in Simulation and Transfers Zero-Shot to Physical Hardware}\par\vspace{6pt}
{\small\textbf{By} MIT Han Lab (PEAR Team)}\par\vspace{2pt}
{\small\color{keywordgray}arXiv:2611.18562 \textbullet{} cs.RO, cs.LG \textbullet{} Submitted 13 Nov 2026}
\par\vspace{2pt}\noindent{\color{rulered}\rule{\columnwidth}{0.6pt}}\vspace{6pt}

A 27-billion-parameter open-weight reasoning model, trained entirely via reinforcement learning in simulation with no expert demonstrations, achieves 74.4\% task success in the LIBERO suite---up from 1.6\%---and transfers zero-shot to a real dual-arm robot at 60\% success (18\,/\,30 trials), where the base model succeeds in none.
The MIT Han Lab's PEAR system, presented in arXiv:2611.18562, demonstrates that the data-collection bottleneck of robot learning can be bypassed when a large reasoning model is paired with scalable simulation-based RL.

\begin{mdframed}[backgroundcolor=boxgray,linecolor=rulered,linewidth=1pt,
    innerleftmargin=6pt,innerrightmargin=6pt,
    innertopmargin=6pt,innerbottommargin=6pt]
\textbf{\small AT A GLANCE}\par\vspace{4pt}
\begin{description}[leftmargin=0pt,labelsep=4pt,itemsep=2pt,parsep=0pt,
    font=\small\bfseries]
  \item[Problem:] Generalist robot policies require large-scale expert demonstration data, which is expensive and embodiment-specific.
  \item[Why it matters:] Eliminating demonstrations removes the main cost and embodiment barrier to deploying robot learning at scale.
  \item[Method:] Robot-use RL (RURL) in simulation with binary success rewards; PEAR orchestrates rollouts via vLLM+Ray and training via veRL.
  \item[Result:] LIBERO: 1.6\%$\rightarrow$74.4\%; real robot (zero-shot): 0/30$\rightarrow$18/30; RoboLab-120: 37.3\%.
  \item[Evidence:] LIBERO, RoboDojo, RoboLab-120 (simulation); YAM dual-arm robot (real hardware, 30 trials).
\end{description}
\end{mdframed}

\section{The Big Picture}

Robot learning has long depended on demonstration data: a human operator teleoperates a robot to collect thousands of successful trajectories, which are then used for imitation learning.
This pipeline works, but it scales poorly: each new embodiment needs its own demonstrations, and rare failure modes are hard to cover.

Reinforcement learning from scratch offers an alternative---let the robot explore and learn from outcomes---but standard RL with physical robots is slow and risky.
Simulation bridges this gap in principle, but prior approaches required dense reward shaping or privileged simulation state that does not exist on real hardware.

PEAR asks whether a large reasoning model can close the loop differently: use the model's pre-existing planning capabilities to bootstrap exploration in simulation, provide only binary task completion rewards, and trust the model's generalization to carry learned behavior to physical hardware.
The answer is a striking affirmative.

\section{Why Existing Approaches Fall Short}

Imitation learning methods such as ACT, Diffusion Policy, and $\pi_0$ achieve strong performance but require demonstration data at every new task.
Scaling them to diverse tasks multiplies data collection cost linearly.

RL-from-demonstrations hybrids reduce data requirements but do not eliminate them; they also inherit the coverage problem---rare object configurations and grasps that do not appear in demonstrations are handled poorly.

Previous sim-to-real transfer work for manipulation typically required dense reward functions (e.g., distance to target, grasp success sub-stages) that must be manually engineered per task.
Engineering these rewards for the breadth of tasks in benchmarks like LIBERO-Plus is a significant research effort in itself.

\section{Core Mechanism}

PEAR trains a 27B reasoning model (Qwen3.8-27B) on robotic tasks using robot-use RL (RURL):
\begin{enumerate}[leftmargin=1.5em,itemsep=2pt]
  \item \textbf{Observation:} At each step, the model receives an image observation and proprioception, tokenized and fed as context.
  \item \textbf{Action:} The model outputs a motor command sequence via autoregressive decoding.
  \item \textbf{Reward:} A binary signal ($r=1$ at successful task completion, $r=0$ otherwise) is provided by the simulator.
  \item \textbf{Update:} PPO updates the model's parameters based on rollout returns.
\end{enumerate}

The two-stage training schedule first trains on LIBERO (225 steps), then extends to RoboDojo for generalization, before evaluation on the held-out RoboLab-120 benchmark.

\section{Technical Formulation}

\subsection*{RURL Objective}

PEAR optimizes the standard RL objective with a sparse reward:
\begin{equation}
\mathcal{J}(\theta) = \mathbb{E}_{\tau \sim \pi_\theta}\!\left[\sum_{t=0}^{T} \gamma^t r_t\right], \quad r_t = \mathbf{1}[\text{task complete at } t]
\label{eq:rurl}
\end{equation}
No dense reward shaping, no auxiliary losses, and no demonstration data enter the objective.
The model must discover manipulation strategies from binary outcome feedback alone.

\subsection*{PPO Clipped Surrogate}

Policy gradient updates use the PPO clipped surrogate:
\begin{equation}
\mathcal{L}^{\text{CLIP}}(\theta) = \mathbb{E}_t\!\left[\min\!\left(r_t(\theta)\hat{A}_t,\; \mathrm{clip}(r_t(\theta),1\pm\epsilon)\hat{A}_t\right)\right]
\label{eq:ppo}
\end{equation}
where $r_t(\theta) = \pi_\theta(a_t \mid s_t) / \pi_{\theta_\text{old}}(a_t \mid s_t)$ is the probability ratio and $\hat{A}_t$ is the generalized advantage estimate computed from the binary reward trajectory.
The clip threshold $\epsilon = 0.2$ is standard.

\section{Learning or Inference Procedure}

\textbf{Training (PEAR infrastructure):}
\begin{itemize}[leftmargin=1em,itemsep=1pt]
  \item \textbf{Rollout:} vLLM + Ray generate episodes in the LIBERO simulator across $N$ parallel workers.
  \item \textbf{KV-cache:} LMCache reuses observation encodings across timesteps within an episode, reducing decoding cost by approximately 40\%.
  \item \textbf{Update:} veRL performs PPO updates, with periodic checkpoint saves.
\end{itemize}

\textbf{Inference (zero-shot transfer):}
The trained checkpoint is loaded into the same 27B model and queried with real sensor observations, with no fine-tuning or domain adaptation.
The model generates motor commands at $\approx$10 Hz on the YAM dual-arm robot.

\section{What the Guarantee Says}

No formal guarantee is provided in the paper.
The empirical evidence for zero-shot sim-to-real transfer rests on 30 trials, which is a limited sample.
The authors note that the binary reward signal provides no gradient signal for failed episodes until the model first achieves success---a cold-start problem that the model's pre-trained planning capabilities appear to resolve in practice, but which could fail for tasks with very low initial success probability.

\section{Experimental Findings}

\begin{center}
\small
\begin{tabular}{lrr}
\toprule
\textbf{Benchmark} & \textbf{Base} & \textbf{PEAR} \\
\midrule
LIBERO (sim)         & 1.6\%  & 74.4\% \\
RoboLab-120 (sim)    & 5.4\%  & 22.1\% \\
RoboLab-120 w/ RoboDojo & —    & 37.3\% \\
Real robot (30 trials) & 0\%  & 60.0\% \\
\bottomrule
\end{tabular}
\end{center}

The base model's near-zero LIBERO success rate reflects a fundamental mismatch between pre-training tasks (language) and the required behavior (motor control).
PEAR's 225-step training on LIBERO is enough to bridge this gap.
The generalization to RoboLab-120 from LIBERO-only training (22.1\%) confirms that the model learns generalizable rather than task-specific behavior.

\section{Ablations and Interpretation}

\begin{itemize}[leftmargin=1em,itemsep=2pt]
  \item \textbf{LMCache impact:} Disabling KV-cache reuse reduces throughput by approximately 40\%, making training impractically slow at the 27B scale.
  \item \textbf{Model size:} Smaller variants (7B, 14B) trained with PEAR reach lower LIBERO plateaus, consistent with scaling laws for reasoning tasks.
  \item \textbf{Reward sparsity:} Adding dense intermediate rewards (joint-velocity penalty, grasp detection) does not improve final performance and sometimes hurts generalization, confirming that binary task rewards are sufficient.
  \item \textbf{Two-stage schedule:} Skipping RoboDojo and training only on LIBERO reduces RoboLab-120 performance from 37.3\% to 22.1\%, showing that domain diversity during RL is important for generalization.
\end{itemize}

\section*{Reference}
MIT Han Lab.
\textbf{Training Reasoning Models in Simulation to Natively Use Real Robots.}
arXiv:2611.18562. cs.RO, cs.LG. November 2026.
\url{https://arxiv.org/abs/2611.18562}

\end{document}
